\section{Padrões Arquiteturais}

Padrões podem ser definidos como conceitos reutilizáveis que podem ser aplicados para resolver questões similares em sistemas em contextos similares, considerando necessidades não funcionais. Isso significa que um padrão pode ser aplicado a qualquer sistema, independentemente de suas funcionalidades, desde que tenham requisitos de qualidade similares.

Se funcionalidades fossem a única parte importante, não haveria necessidade de ter arquitetura de software. Mas ao mesmo tempo, isso não é o que acontece na realidade. Requisitos de qualidade acabam não sendo tão relevantes quanto deveriam ao escolher padrões arquiteturais para um sistema na indústria, e isso traz consequências de longo prazo para confiabilidade, manutenibilidade, testabilidade, entre outras complicações.
